% Перерисовка исторической схемы Vallado по исходному рисунку лекции 11.
% Положения кривых восстановлены графически, обе шкалы логарифмические.
\begin{tikzpicture}[course diagram,x=1mm,y=.95mm]
\path[use as bounding box] (-11,-7) rectangle (85,53);
\node[anchor=south] at (36,51) {Погрешность положения};
\foreach \y/\lab in {0/{10 см},10/{1 м},20/{10 м},30/{100 м},40/{1 км},50/{10 км}}{
 \draw[aux line] (0,\y)--(75,\y);
 \node[anchor=east,inner sep=1mm] at (0,\y) {\lab};
}
\foreach \x/\lab in {0/{100},25/{1000},50/{10\,000},75/{100\,000}}{
 \draw[aux line] (\x,0)--(\x,50);
 \node[below,inner sep=1mm] at (\x,0) {$\lab$};
}
% Промежуточные деления логарифмических шкал.
\foreach \d in {0,1,2}{\foreach \k in {2,...,9}{
 \pgfmathsetmacro{\tickx}{25*(\d+ln(\k)/ln(10))}
 \draw[coursemuted,line width=.35pt] (\tickx,0)--(\tickx,-.7);
}}
\foreach \d in {0,...,4}{\foreach \k in {2,...,9}{
 \pgfmathsetmacro{\ticky}{10*(\d+ln(\k)/ln(10))}
 \draw[coursemuted,line width=.35pt] (0,\ticky)--(-.6,\ticky);
}}
\draw[courseink] (0,50)--(0,0)--(75,0);
\node[anchor=north,inner sep=0pt] at (37.5,-4.5) {$H$, км};
% 1--3: аналитический прогноз, аппроксимация наблюдений, теория.
\draw[courseaccent,line width=1.1pt]
 (3.7,50) .. controls (5.5,45) and (18,43.7) .. (32.5,43.7)
 .. controls (45,43.7) and (58,45.2) .. (64,49.8);
\draw[courseaccent,line width=.65pt]
 (3.8,45.2) .. controls (8,41.6) and (20,40) .. (35,40)
 .. controls (50,40) and (63,41.5) .. (68.4,45.4);
\draw[courseaccent,line width=1.3pt] (2.5,39)--(71,39);
% 4--5: РЛС низкой точности и интерферометр.
\draw[coursemuted,dashed,line width=.65pt]
 (3,34.9) .. controls (34,34.5) and (50,34.4) .. (72.5,36.5);
\draw[coursemuted,dashed,line width=.65pt] (2.5,33)--(45,33);
% 6--9: численный прогноз, РЛС высокой точности, теория, лазерные измерения.
\draw[courseink,line width=1.1pt]
 (4.5,21.1) .. controls (10,19.6) and (20,19.4) .. (32,19.4)
 .. controls (47,19.4) and (62,21.1) .. (72.5,23);
\draw[coursemuted,dashed,line width=.65pt]
 (3.4,19.1) .. controls (31,18.4) and (49,18.8) .. (73,21.3);
\draw[courseink,line width=1.3pt] (3,16)--(70,16);
\draw[coursemuted,dashed,line width=.65pt] (3,7.8)--(67.5,7.8);
% Номера размещены с фиксированными промежутками; тонкие линии --- выноски.
\tikzset{accuracy leader/.style={draw=coursemuted,line width=.4pt}}
\draw[accuracy leader] (32.5,43.7)--(32.5,45.4);
\node[anchor=south,inner sep=.3mm,text=courseaccent] at (32.5,45.4) {1};
\draw[accuracy leader] (68.4,45.4)--(72,47);
\node[anchor=south,inner sep=.3mm,text=courseaccent] at (72,47) {2};
\draw[accuracy leader] (45,33)--(48,31);
\node[anchor=west,inner sep=.6mm] at (48,31) {5};
\foreach \xe/\ye/\yl/\num in {71/39/40/3,72.5/36.5/35/4,72.5/23/25/6,73/21.3/20/7,70/16/15/8,67.5/7.8/7.8/9}{
 \draw[accuracy leader] (\xe,\ye)--(79,\yl)--(80,\yl);
 \node[anchor=west,inner sep=.6mm] at (80,\yl) {\num};
}
\end{tikzpicture}
